% year: תשפ"ג
% semester: ב
% moed: א
% source: exams/מבחנים/תשפג/מועד א תשפג סמסטר ב.pdf
% number: 4ב
% points: 5
% categories: derivatives

\begin{question}
עבור אילו ערכים של $a$ הפונקציה $f(x)=\min\left(x^2+x,a\right)$ תהיה גזירה לכל $x$ ממשי? נמקו את תשובתכם.
\end{question}

\begin{hint}
השוו את $a$ לערך המינימלי של $x^2+x$. אם הישר $y=a$ חותך את הפרבולה, מה קורה לנגזרות החד-צדדיות בנקודת החיתוך?
\end{hint}

\begin{solution}
הפרבולה $p(x)=x^2+x=\left(x+\frac12\right)^2-\frac14$ מקבלת ערך מינימלי $-\frac14$ ב-$x=-\frac12$.
\begin{itemize}
  \item אם $a\le-\frac14$: $p(x)\ge-\frac14\ge a$ לכל $x$, ולכן $f(x)\equiv a$ קבועה, וגזירה בכל $\R$.
  \item אם $a>-\frac14$: למשוואה $x^2+x=a$ שני שורשים $x_1<x_2$, $x_{1,2}=\frac{-1\mp\sqrt{1+4a}}{2}$. אז $f(x)=x^2+x$ על $[x_1,x_2]$ ו-$f(x)=a$ מחוצה לו. בנקודה $x_1$ (שבה $f$ רציפה) הנגזרת השמאלית היא $0$ (פונקציה קבועה), והימנית היא $p'(x_1)=2x_1+1=-\sqrt{1+4a}\neq0$. הנגזרות החד-צדדיות שונות ולכן $f$ אינה גזירה ב-$x_1$ (ובאופן דומה ב-$x_2$, שם הנגזרת השמאלית $\sqrt{1+4a}$ והימנית $0$).
\end{itemize}
\textbf{תשובה:} $f$ גזירה לכל $x$ ממשי אם ורק אם $\boxed{a\le-\frac14}$.
\end{solution}
